% AgriTwin — IEEE-style draft. Compile with: pdflatex agritwin.tex && bibtex agritwin
% Fill AUTHOR BLOCK and VENUE before submission. All result values are macros at the top.
\documentclass[journal,twocolumn,10pt]{IEEEtran}

\usepackage{amsmath,amssymb}
\usepackage{graphicx}
\usepackage{booktabs}
\usepackage{multirow}
\usepackage{url}
\usepackage{caption}
\usepackage{subcaption}
\usepackage{xcolor}

% ---------------------------------------------------------------------------
% AUTHOR BLOCK — name/institution taken from the project review records;
% fill CITY, COUNTRY and EMAIL before submission
% ---------------------------------------------------------------------------
\author{
  \IEEEauthorblockN{Jashwanth Aravapalli}\\
  \IEEEauthorblockA{Department of CSE (AI \& ML), Keshav Memorial Engineering College\\
  [CITY], [COUNTRY]\\
  [EMAIL]}
}

% ---------------------------------------------------------------------------
% MEASURED RESULT MACROS — update here only when new runs finish
% ---------------------------------------------------------------------------
\newcommand{\PVvtwo}{0.9511}
\newcommand{\PDvtwo}{0.6780}
\newcommand{\FonePVv2}{0.9507}
\newcommand{\FonePDv2}{0.6501}
\newcommand{\EonePV}{0.9549}
\newcommand{\EonePD}{0.3220}
\newcommand{\EtwoPV}{0.3120}
\newcommand{\EtwoPD}{0.6822}
\newcommand{\EthreePDcontrolled}{0.6653}
\newcommand{\EfourPV}{0.8308}
\newcommand{\EfourPD}{0.5551}
\newcommand{\ExGLab}{0.9586}
\newcommand{\ExGField}{0.6483}
\newcommand{\RegLab}{0.9549}
\newcommand{\RegField}{0.6610}
% Deployed v3 — ConvNeXt-Tiny with metadata-driven 4-view TTA
\newcommand{\PVvthreeSingle}{0.9774}
\newcommand{\PDvthreeSingle}{0.6907}
\newcommand{\FonePDvthreeSingle}{0.6833}
\newcommand{\PVvthree}{0.9812}
\newcommand{\PDvthree}{0.7161}
\newcommand{\FonePDvthree}{0.7093}
\newcommand{\PrecPDvthree}{0.7373}
\newcommand{\RecPDvthree}{0.7106}
\newcommand{\DeltaTTA}{0.0254}
\newcommand{\LowConfSingle}{0.3333}
\newcommand{\LowConfTTA}{0.4167}
% Soft crop prior: external farm crop hint, alpha = 0.95
\newcommand{\CropAlpha}{0.95}
\newcommand{\PDprior}{0.8093}
\newcommand{\PDpriorWrong}{0.7754}
\newcommand{\PDhard}{0.8305}
\newcommand{\PDhardWrong}{0.6568}

\begin{document}

\title{AgriTwin: Uncertainty-Aware Fusion of Leaf, Yield and Market Signals for Crop Decisions,
with Cross-Domain Validation of the Disease Input}

\maketitle

\begin{abstract}
Farmers combine a leaf photograph, a yield expectation, a mandi price, a buyer offer and a transport
quote to decide what to grow, when to harvest and where to sell. These signals are usually processed
in isolation, so an individually accurate prediction need not produce a good economic decision.
We present \emph{AgriTwin}, an agricultural digital-twin platform that unifies crop-disease
screening, conformal yield prediction, market forecasting and transport costing in one
uncertainty-bearing state consumed by a deterministic decision engine. Because the leaf photo is
the first link in this chain, a confidently wrong disease prediction would corrupt the entire
decision; we therefore validate that input under domain shift. A controlled experiment shows that
an EfficientNet-B0 trained only on PlantVillage reaches \EonePV accuracy in the lab but collapses
to \EonePD on in-the-wild PlantDoc photographs, a 63.3-point gap. Joint training retains lab
performance (\PVvtwo) and more than doubles field accuracy (\EthreePDcontrolled in the
controlled run; \PDvtwo for the independently retrained v2 checkpoint). A follow-up ConvNeXt-Tiny
trained on a de-duplicated joint manifest is now deployed with metadata-driven 4-view TTA, reaching
\PVvthree lab / \PDvthree field accuracy (field macro-F1 \FonePDvthree); a soft crop prior from
the selected farm crop reaches \PDprior field accuracy when the hint is correct. Contrary to the usual
assumption, sequential fine-tuning underperformed joint training on both domains
(\EfourPV lab / \EfourPD field). Grad-CAM overlays confirm lesion-focused attention. The system
turns these predictions, yield and price intervals into ranked sell/wait decisions with an
explicit INSPECT fallback for low-confidence inputs. All numbers are reproduced by released code
and stored in artefact metadata.
\end{abstract}

\begin{IEEEkeywords}
agricultural digital twin, decision support, plant disease detection, domain adaptation,
uncertainty quantification, Grad-CAM, EfficientNet, ConvNeXt, test-time augmentation
\end{IEEEkeywords}

\section{Introduction}
\label{sec:intro}

Farmers' economically important decisions are multi-source. Crop health is judged from a leaf
photo, expected production from experience or a model, the selling price from yesterday's mandi
quote, a buyer may offer a different price by phone, and transport cost depends on distance and
vehicle. Each source is uncertain, time-stamped and perishable. Treating them as separate
predictions leaves the farmer to perform an implicit, untracked fusion; an individually good
prediction does not guarantee a good decision.

AgriTwin represents all sources as one typed, uncertainty-bearing decision state. Disease
screening returns a label, a confidence and a Grad-CAM explanation; yield prediction returns a
point estimate and an 80\% split-conformal interval; market forecasting returns a price and
empirical residual bounds; buyer offers carry validity windows; and transport returns a cost
breakdown. A pure decision engine combines these into eight weighted risk components, two revenue
scenarios (sell now / wait) and an ordered rule set recommending SELL, WAIT or INSPECT, with a
human-readable explanation.

This design makes the reliability of the first signal a system-level concern. Deep learning
studies routinely report 95--99\% accuracy on PlantVillage, where single leaves are photographed
face-up against uniform backgrounds. Field photographs contain variable illumination, occlusion,
multiple leaves, soil and sky. A model that exploits the controlled background as a shortcut will
be confidently wrong in exactly the conditions the decision system is meant to support.

The contributions are:
\begin{enumerate}
  \item an integrated uncertainty-aware decision system fusing leaf health, yield intervals,
        market forecasts, buyer offers and transport cost into one reproducible decision with
        explanations and an INSPECT fallback;
  \item a measured quantification of the lab-to-field failure (\EonePV $\to$ \EonePD) using the
        official PlantDoc test split, with no test image used for training or selection;
  \item a controlled four-protocol comparison showing joint training outperforms sequential
        fine-tuning, including the honest negative result that two-stage adaptation loses on
        both domains;
  \item explainability and efficiency evidence: Grad-CAM verification and a measured
        parameter/FLOP/latency profile supporting near-real-time inference.
\end{enumerate}

\section{Related Work}

\emph{Plant disease detection.} Mohanty \emph{et al.}~\cite{mohanty2016} demonstrated deep
convolutional networks on PlantVillage and already noted its limited realism. Transfer learning
with EfficientNet, ResNet and MobileNet dominates subsequent work, usually with in-domain
accuracy only; surveys document the gap to field readiness~\cite{survey2025,review2025}.

\emph{In-the-wild data.} PlantDoc~\cite{plantdoc} provides internet photographs with natural
backgrounds and multiple leaves. Recent lightweight-model evaluations report in-domain PlantDoc
results and edge deployment~\cite{edge2026}.

\emph{Cross-domain detection.} Deployable cross-domain smartphone systems exist~\cite{crossdomain2025}.
Our classifier is instead treated as an input-reliability component of a larger decision system,
and we report the negative two-stage result rather than only the best configuration.

\emph{Explainability and digital twins.} Grad-CAM~\cite{gradcam} localises evidence without
architecture changes and is used here as a qualitative check and a user-facing output; reviews
recommend it to detect background shortcuts~\cite{review2025}. Agricultural digital twins have
demonstrated value for individualised management~\cite{kim2024}. We work at their intersection:
a decision pipeline whose first input is cross-domain validated and explained.

\emph{Excess-green prior.} The excess-green index highlights vegetation and suppresses soil. An
ExG attention gate was effective in our complementary UAV crop/weed study (ReVit). Here the same
prior is adapted as a training-time background-suppression ablation, reported only if it improves
field accuracy; the two systems remain separate and cross-cited.

\section{System and Methods}
\label{sec:methods}

\subsection{System architecture}

The implemented system has seven layers: data; persistence (SQLAlchemy 2, Alembic); ML inference
(ConvNeXt-Tiny with 4-view TTA + Grad-CAM, ExtraTrees + conformal intervals, market forecasters); a pure decision
engine (policy $\to$ risk $\to$ fusion $\to$ ordered rules); service orchestration; a FastAPI
layer with JWT, validation and rate limiting; and a React 18/TypeScript presentation layer.

A scan is validated (MIME allow-list, 8\,MB and 4096\,px limits, SHA-256 deduplication), passed to
the crop-health service, and persisted with its Grad-CAM overlay. The decision service reads this
prediction, the latest yield prediction, market forecast, buyer offers and transport estimate,
builds typed inputs, calls the engine, and stores the decision and factors. Every response labels
its model mode (\texttt{real} vs \texttt{demo}), so fallback behaviour is never silent.

\subsection{Datasets}

\begin{table}[t]
\centering
\caption{Datasets and split sizes.}
\label{tab:data}
\begin{tabular}{@{}llr@{}}
\toprule
Dataset & Role & Images \\
\midrule
PlantVillage subset & controlled lab, disease semantics & 1{,}900 \\
PlantDoc (CC BY 4.0) & in-the-wild internet/field images & 2{,}572 \\
Web / farm photos & optional test-only external domain & user-supplied \\
\bottomrule
\end{tabular}
\end{table}

PlantVillage contributes 1{,}368/266/266 train/validation/test images; PlantDoc contributes
1{,}998/338/236. A deterministic manifest stores all 4{,}472 entries with a domain tag. Of 2{,}578
PlantDoc files, 2{,}572 were recovered; six filenames with Windows-illegal characters could not be
stored, while all 28 classes and both official splits remain complete. The official PlantDoc test
split is never used for training or early stopping. Twenty-eight PlantDoc folders map onto the
same fixed 38-class output space used by the deployed service. A SHA-256 audit found 12 duplicate
image groups, 11 crossing train/validation and test splits; the de-duplicated manifest used by v3
contains 4{,}460 entries and preserves every held-out test image.

\subsection{Model and preprocessing}

The controlled study used ImageNet-pretrained EfficientNet-B0 with a 38-logit head; the deployed v3
model uses ConvNeXt-Tiny with the same fixed class order. Inputs are 192$\times$192 central crops
normalised with ImageNet statistics. Training augmentation: RandomResizedCrop (scale 0.65--1.0),
horizontal flip, vertical flip ($p=0.15$), rotation up to $15^\circ$, colour jitter and random
grayscale ($p=0.05$). Evaluation uses deterministic resize and centre crop; deployed serving
averages four views (original and horizontal flip at scales 1.0 and 0.9) before argmax. Measured
complexity on an RTX 5060 Laptop with CUDA 12.8 wheels: ConvNeXt-Tiny has 27{,}849{,}350 parameters,
a 106.24\,MB checkpoint and 6.5458 GFLOPs, with 31.9\,ms/image at batch 1 and about 231 images/s at
batch 32 for one forward pass. The historical EfficientNet-B0 run had 4{,}056{,}226 parameters and
0.5655 GFLOPs. Full data: \texttt{benchmark\_results\_v3.json}.

\subsection{Training and cross-domain protocol}

Inverse-frequency weighted cross-entropy, AdamW (3e-4 for the historical study, 1e-4 for v3,
weight decay 1e-4), cosine schedule and early stopping on validation macro-F1; seed 42; FP16
autocast with GradScaler on CUDA. Optional
label smoothing and mixup are metadata-recorded flags. Four protocols run from the same manifest:
\texttt{baseline\_pv} (lab only), \texttt{plantdoc\_only} (field only), \texttt{mixed} (joint) and
\texttt{two\_stage} (lab then field). Each is evaluated on both held-out test sets.

\subsection{Explainability and decision fusion}

Grad-CAM~\cite{gradcam} weights the last convolutional feature maps (\texttt{features[-1]} for
EfficientNet, \texttt{features[-1][-1]} for ConvNeXt) by globally averaged gradients
of the predicted class score, applies ReLU and upsamples. Overlays are produced offline for the
paper and online in the API. The decision engine weights eight risk components (yield uncertainty,
crop health, offer freshness/competitiveness, market uncertainty, transport, completeness, price
staleness, deterioration). Health confidence gates the INSPECT rule at 0.55: below it, the health
signal is discounted and manual verification may replace an economic action. Yield and price
intervals propagate into revenue bounds.

\section{Results}
\label{sec:results}

\subsection{Main cross-domain comparison}

\begin{table}[t]
\centering
\caption{Accuracy and macro-F1 on held-out lab (PV) and field (PD) test sets.}
\label{tab:cross}
\begin{tabular}{@{}lcccc@{}}
\toprule
Protocol & PV acc & PD acc & PV F1 & PD F1 \\
\midrule
\texttt{baseline\_pv}   & \textbf{\EonePV} & \EonePD & 0.9539 & 0.2177 \\
\texttt{plantdoc\_only} & \EtwoPV & \EtwoPD & 0.2396 & 0.6721 \\
\texttt{mixed}          & \PVvtwo & \textbf{\EthreePDcontrolled} & \textbf{\FonePVv2} & 0.6563 \\
\texttt{two\_stage}     & \EfourPV & \EfourPD & 0.8322 & 0.5292 \\
\bottomrule
\end{tabular}
\end{table}

The lab-only model loses 63.3 accuracy points on real photographs. Joint training retains lab
accuracy (\EonePV $\to$ \PVvtwo) while more than doubling field accuracy (\EonePD $\to$
\EthreePDcontrolled). The expected winner, two-stage fine-tuning, underperforms joint training on
\emph{both} domains, consistent with catastrophic forgetting of the lab domain. The deployed
checkpoint was independently retrained to save the winning weights and regenerate Grad-CAM; it
measured \PVvtwo lab / \PDvtwo field (F1 \FonePVv2 / \FonePDv2). The small field difference
from the controlled run is attributed to non-deterministic CUDA kernel selection; both are stored
in the v2 backup metadata.

\subsection{Deployed ConvNeXt-Tiny and test-time augmentation}

\begin{table}[t]
\centering
\caption{Deployed v3 on held-out tests: single view versus 4-view TTA.}
\label{tab:v3}
\begin{tabular}{@{}lccc@{}}
\toprule
Inference & PV acc & PD acc & PD F1 \\
\midrule
Single view & \PVvthreeSingle & \PDvthreeSingle & \FonePDvthreeSingle \\
4-view TTA (deployed) & \textbf{\PVvthree} & \textbf{\PDvthree} & \textbf{\FonePDvthree} \\
\bottomrule
\end{tabular}
\end{table}

The ConvNeXt-Tiny candidate was promoted only because its single-view field accuracy \PDvthreeSingle
strictly exceeded the v2 baseline \PDvtwo. Metadata-driven 4-view TTA was then enabled after a
further \DeltaTTA gain on the same held-out PlantDoc split: field macro precision \PrecPDvthree and
recall \RecPDvthree, nine images fixed and three broken. Below the 0.55 screening threshold,
accuracy rose from \LowConfSingle to \LowConfTTA ($n=48$). Serving and \texttt{tta\_eval.py} call
the same transform helper, and CUDA and CPU evaluations both return \PDvthree.

The scan form already requests a farm crop, so we also measured an external soft crop prior
(\texttt{alpha=\CropAlpha}) that reweights probabilities toward classes of the selected crop before
argmax. On the same held-out PlantDoc split it reaches \PDprior with a correct hint and
\PDpriorWrong when 20\% of hints are wrong; unsupported crops fall back unchanged. A hard mask
(\texttt{alpha=1}) reaches \PDhard with correct hints but collapses to \PDhardWrong at 20\% wrong
hints, so it was rejected. The soft prior is a serving enhancement measured with
\texttt{crop\_prior\_eval.py}; weight promotion still rests on \PDvthreeSingle / \PDvthree.

\subsection{Measured ablations}

\begin{table}[t]
\centering
\caption{Additional ablations. TTA is model-dependent: rejected for v2, enabled for v3.}
\label{tab:ablate}
\begin{tabular}{@{}lccc@{}}
\toprule
Variant & Lab & Field & Outcome \\
\midrule
256\,px, 20 epochs, PD-val selection & 0.9586 & 0.6525 & rejected \\
4-view TTA on v2 EfficientNet-B0 & 0.9662 & 0.6864 & rejected \\
ExG background suppression & \ExGLab & \ExGField & rejected \\
Label smoothing + mixup & \RegLab & \RegField & rejected \\
ConvNeXt-Tiny, de-duplicated manifest & \PVvthreeSingle & \PDvthreeSingle & promoted \\
4-view TTA on v3 ConvNeXt-Tiny & \PVvthree & \PDvthree & enabled \\
\bottomrule
\end{tabular}
\end{table}

More pixels and an excess-green prior are not automatic remedies under domain shift: 256\,px with
field-domain checkpoint selection reduced field accuracy, and ExG suppresses the brown and necrotic
tissue that carries lesion evidence, so it remains in the separate ReVit study and is only
cross-cited here. Label smoothing ($0.1$) with mixup ($0.2$) retained lab accuracy ($0.9549$) but
reduced field accuracy from \PDvtwo to \RegField and lowered macro-F1 from \FonePDvtwo to $0.5584$;
the v2 checkpoint was therefore left unchanged. Test-time augmentation is model-dependent: it was
rejected for v2 (two net images, no low-confidence gain) but enabled for v3, where it added
\DeltaTTA field accuracy and improved low-confidence accuracy from \LowConfSingle to \LowConfTTA at
roughly four times the forward-pass cost.

\subsection{Per-class behaviour and explainability}

Of 38 classes, 27 have non-zero PlantDoc test support. For the promoted single-view ConvNeXt
checkpoint, the weakest supported classes are maize gray leaf spot (F1 0.143, $n=4$), tomato
bacterial spot (0.235, $n=9$), soybean healthy (0.364, $n=8$) and potato early blight (0.462,
$n=8$). Dominant confusions are biologically plausible: maize northern leaf blight $\to$ maize gray
leaf spot (7), tomato bacterial spot $\to$ tomato septoria leaf spot (6), pepper bacterial spot
$\to$ tomato bacterial spot (4), peach healthy $\to$ cherry healthy (4) and potato early blight
$\to$ potato late blight (3). Grad-CAM overlays generated from the deployed model are referenced in
metadata and show attention on lesion regions rather than uniform background in correctly
classified examples.

\subsection{Efficiency and system behaviour}

\begin{table}[t]
\centering
\caption{Measured efficiency, RTX 5060 Laptop, 192\,px input, CUDA 12.8 wheels. One forward pass.}
\label{tab:eff}
\begin{tabular}{@{}lrrrr@{}}
\toprule
Architecture & Params & Size & GFLOPs & ms/img \\
\midrule
MobileNet-V3-Large & 4.25M & 16.22\,MB & 0.3179 & 33.3 \\
EfficientNet-B0 (historical v2) & 4.06M & 15.47\,MB & 0.5655 & 42.3 \\
EfficientNet-B3 & 10.75M & 41.03\,MB & 1.4134 & 66.7 \\
EfficientNet-V2-S & 20.23M & 77.16\,MB & 4.1878 & 94.9 \\
ResNet-50 & 23.59M & 89.97\,MB & 6.0057 & 31.5 \\
ConvNeXt-Tiny (deployed v3) & 27.85M & 106.24\,MB & 6.5458 & 31.9 \\
ConvNeXt-Small & 49.48M & 188.77\,MB & 12.7587 & 54.5 \\
\bottomrule
\end{tabular}
\end{table}

Source: \texttt{benchmark\_results\_v3.json}; the earlier cu118 run is preserved in
\texttt{benchmark\_results.json}. Deployed v3 averages four views, about 1.76\,s per CPU
prediction in the smoke test. A helper runs arbitrary user photos through the service and reports
label, confidence, severity and ACCEPT/INSPECT verdict; on five downloaded photos it returned 2/5
exact matches and routed two low-confidence cases to INSPECT. These are smoke checks, not a labelled
test set; see \texttt{docs/CROP\_HEALTH\_V3.md}.

\section{Discussion}

The 63-point collapse is consistent with shortcut learning: a uniform background and consistent
leaf orientation let a model score highly without learning robust disease features. Joint training
works because field images break the shortcut while lab images preserve clean disease semantics;
sequential fine-tuning instead overwrites the lab representation.

The system-level implication is that confidence must be first-class. A classifier confidently
wrong under domain shift is more dangerous than an uncertain one. The INSPECT rule, confidence
discounting and Grad-CAM overlay are therefore the mechanism by which an imperfect input is kept
from dominating an economic decision.

Limitations remain: PlantDoc has roughly 2.6k images with single-label annotations of possibly
multi-disease leaves; the model assumes one leaf per image and covers 14 crop species; abiotic
stress can mimic disease; and the optional web test set was unavailable for this draft. Field
validation with longitudinal farmer data, multi-label severity scoring and active learning are
needed before the screening output can be treated as diagnostic.

\section{Conclusion and Future Work}

We presented AgriTwin, an uncertainty-aware agricultural decision system fusing disease, yield,
market, offer and transport inputs into a single explainable decision. Its leaf-photo input was
cross-domain validated: the lab-to-field collapse was measured, joint training recovered field
accuracy while retaining lab accuracy, and sequential adaptation failed. The deployed ConvNeXt-Tiny
model reaches \PVvthree lab / \PDvthree field accuracy with 4-view TTA, or \PDprior with a correct
crop hint, while remaining a screening aid; the engine falls back to INSPECT when confidence is low.
Future work includes
on-device deployment, external field trials, multi-label disease and severity modelling, and
active learning from farmer feedback.

\section*{Acknowledgment}
[ACKNOWLEDGMENTS / FUNDING, IF APPLICABLE]

\bibliographystyle{IEEEtran}
\bibliography{references}

\end{document}
